\chapter{Global Types}\label{ch:global}

A LeanFM protocol is a \emph{global type}: one description of a conversation among
named roles, written from no participant's point of view. This chapter introduces the
language through two protocols that recur throughout the book. Every later chapter
reduces what is written here to CTL, so each construct is introduced with the question
it raises for the reduction.

\section{A request with two answers}

The smallest interesting protocol is a request with more than one possible answer.
A browser asks for an image; the server returns it or says it is not there.

\begin{lstlisting}[language={},basicstyle=\ttfamily\footnotesize,breaklines=true]
structure GetImage where
  id : String

structure Image where
  id : String
  bytes : String

protocol Images (Browser, Server) {
  g: GetImage from Browser to Server route "GET /images/{id}";
  choice {
    i: Image from Server to Browser cause g status 200 where i.id = g.id;
  } or {
    NotFound from Server to Browser cause g status 404;
  }
}
\end{lstlisting}


\begin{figure}[ht]
\centering
\AccessibleFigure{Interaction diagram of Images. Message from Browser to Server: GetImage. Then one of two alternatives. Message from Server to Browser: Image status 200. Or message from Server to Browser: NotFound status 404.}{\ThemedDiagram{width=.58\textwidth}{ctl-images}}
\caption{The Images protocol as LeanFM draws it (\texttt{session\_diagram}).}
\label{fig:images}
\end{figure}

Read it line by line.

\begin{itemize}
\item \texttt{protocol Images (Browser, Server)} declares the roles. A role is a
  part in the conversation, not a machine; one process may play several roles, and
  many processes may play the same role in different sessions.
\item \texttt{g: GetImage from Browser to Server} is a \emph{message}: a label, a
  sender and a receiver. The label names a structure declared above it, so the message
  carries an \texttt{id}. The prefix \texttt{g:} names this occurrence so later
  messages can refer to it.
\item \texttt{route "GET /images/\{id\}"} and \texttt{status 200} say how the message
  travels when the conversation is HTTP. They do not change the order of the
  conversation; they change the interface LeanFM generates.
\item \texttt{media "image/png"} says the answer's body is the image itself, not JSON.
  Image's fields are all declared \texttt{Header}, so they travel as headers beside the
  bytes; LeanFM refuses a field that would have to be in the body. A header is named
  from its field: \texttt{cacheControl} is \texttt{Cache-Control}, \texttt{contentLength}
  is \texttt{Content-Length}.
\item The three lines without a type are \emph{constraints}: propositions every Image
  satisfies. \texttt{cached} fixes the values \texttt{Cache-Control} may take,
  \texttt{small} bounds the image's \texttt{size} in bytes, and \texttt{measured} says the
  \texttt{Content-Length} header tells the truth about it.
\item \texttt{cause g} says the answer answers \emph{that} request. A cause is an
  edge between two events of one run, not between two labels; it is what lets a log
  with many interleaved requests be read back as separate conversations.
\item \texttt{choice \{ \ldots\ \} or \{ \ldots\ \}} offers alternatives. The role that
  sends the first message of a branch makes the choice, here the server.
\item \texttt{where i.id = g.id} is a condition on the data: the image returned is the
  one asked for. Chapter~\ref{ch:where} gives where clauses their own chapter.
\end{itemize}

The question this protocol raises for the reduction is the one in the cover formula:
after the request, \emph{both} answers must be possible, \emph{one} of them must come
next, and \emph{nothing else} may. Chapter~\ref{ch:faithful} turns exactly that sentence
into CTL.

\section{It is already an HTTP API}\label{sec:openapi}

The \texttt{route} and \texttt{status} lines make Images an HTTP interface, and LeanFM
writes it in the standard everyone already reads: an OpenAPI~3.1 (Swagger) document
for the Server, the role that receives requests.

\begin{lstlisting}[language={},basicstyle=\ttfamily\footnotesize,breaklines=true]
openapi: "3.1.0"
info:
  description: "The requests Server receives in protocol Images, generated by LeanFM from the session type."
  title: "Images: Server"
  version: "0.1.0"
paths:
  "/images/{id}":
    get:
      operationId: "GetImage"
      parameters:
        - in: "path"
          name: "id"
          required: true
          schema:
            type: "string"
      responses:
        "200":
          content:
            "application/json":
              schema:
                "$ref": "#/components/schemas/Image"
          description: "Image"
        "404":
          content:
            "application/json":
              schema:
                "$ref": "#/components/schemas/NotFound"
          description: "NotFound"
      summary: "GetImage from Browser"
components:
  schemas:
    GetImage:
      additionalProperties: false
      properties:
        id:
          type: "string"
      required:
        - "id"
      title: "GetImage"
      type: "object"
    Image:
      additionalProperties: false
      properties:
        bytes:
          type: "string"
        id:
          type: "string"
      required:
        - "id"
        - "bytes"
      title: "Image"
      type: "object"
    NotFound:
      additionalProperties: false
      properties: {}
      required: []
      title: "NotFound"
      type: "object"
\end{lstlisting}


Each request is an operation under its route; its answers are responses under their
statuses. The Image is \texttt{image/png} bytes of at most 5{,}000{,}000 bytes, with
its headers and their allowed values, because the structure says so; a message without
a media type is JSON whose schema admits exactly its fields, as NotFound is. What a
schema cannot say, such as \texttt{measured}, LeanFM still checks on every log. Here a
server sent a \texttt{Content-Length} that does not match the bytes:

\begin{lstlisting}[language={},basicstyle=\ttfamily\footnotesize,breaklines=true]
{"fields":{"id":"cat"},"from":"Browser","id":"a1","label":"GetImage","to":"Server"}
{"cause":"a1","fields":{"cacheControl":"no-store","contentLength":"48213","imageId":"cat","size":"48000"},"from":"Server","id":"a2","label":"Image","to":"Browser"}
\end{lstlisting}


\begin{lstlisting}[language={},basicstyle=\ttfamily\footnotesize,breaklines=true]
Images: 1 problem with the logged values (constraints, links, session keys, versions); first: session chain from a1, event a2: Image.measured does not hold: contentLength = size (contentLength = 48213, size = 48000); no Image can be made with these fields
session chain from a1, event a2: Image.measured does not hold: contentLength = size (contentLength = 48213, size = 48000); no Image can be made with these fields
\end{lstlisting}

The \texttt{\{id\}} in the route is a path parameter. Load the document into Swagger UI
or any OpenAPI generator and you have the server's stubs and a client, in whatever
language those tools support. The point of this book is what the OpenAPI document
\emph{cannot} say: that every GetImage is answered, by exactly one of these two
responses, and nothing else. The protocol says that, and Part~\ref{part:ctl} turns it
into CTL that can be checked. The OpenAPI document is the protocol's interface; the
protocol is the interface plus the order.

\section{Three roles, nested choices, a secret}

A purchase involves a client, a shop and a bank, and it can go wrong in two places.

\begin{lstlisting}[language={},basicstyle=\ttfamily\footnotesize,breaklines=true]
import LeanFM.Session.Syntax

/-- What the client orders. -/
structure CreateOrder where
  item : String
  quantity : Nat

/-- The shop's price for the order, in cents. -/
structure Quote where
  price : Nat

/-- The client pays; the card number is a secret, never logged. -/
structure Pay where
  amount : Nat
  card_secr : String

structure Charge where
  amount : Nat
  card_secr : String

structure Receipt where
  amount : Nat

/-- Placing an order: quote, pay, charge the card. -/
protocol PlaceOrder (Client, Shop, Bank) {
  o: CreateOrder from Client to Shop route "POST /orders";
  choice {
    q: Quote from Shop to Client cause o status 201;
    p: Pay from Client to Shop cause q where p.amount = q.price;
    c: Charge from Shop to Bank cause p;
    choice {
      d: Charged from Bank to Shop cause c;
      Receipt from Shop to Client cause d status 200;
    } or {
      x: Declined from Bank to Shop cause c;
      PaymentDeclined from Shop to Client cause x status 402;
    }
  } or {
    SoldOut from Shop to Client cause o status 409;
  }
}
\end{lstlisting}


\begin{figure}[ht]
\centering
\AccessibleFigure{Interaction diagram of PlaceOrder. Message from Client to Shop: CreateOrder. Then either: message from Shop to Client: Quote; message from Client to Shop: Pay; message from Shop to Bank: Charge; then either message from Bank to Shop: Charged and message from Shop to Client: Receipt or message from Bank to Shop: Declined and message from Shop to Client: PaymentDeclined. Or: message from Shop to Client: SoldOut.}{\ThemedDiagram{width=.8\textwidth}{ctl-placeorder}}
\caption{PlaceOrder: a choice nested inside a choice, and a payment whose amount must match the quote.}
\label{fig:placeorder}
\end{figure}

Three things are new. A choice can sit inside a branch of another choice. A field
whose name ends in \texttt{\_secr} is a \emph{secret}: LeanFM can ask where secrets
travel (the \texttt{secret} atom of Chapter~\ref{ch:checker}), and a project decides
whether to forbid, for example, that a card number ever reaches a log. And the where
clause \texttt{p.amount = q.price} relates two messages of the same run: what the
client pays is what the shop quoted.

\paragraph{Not every message is HTTP.} The shop's charge to the bank has no route:
it may be a queue, a socket or a call inside one process. For such messages LeanFM
writes the message types as protobuf (\texttt{session\_types}): one message per
structure, and an \texttt{Event} envelope that carries any of them with the fields
every log needs (who sent it, to whom, and the event it answers). So the same
protocol drives an HTTP service, a gRPC or raw-socket one, and the log both are
checked against.

\begin{lstlisting}[language={},basicstyle=\ttfamily\footnotesize,breaklines=true]
// Generated by LeanFM from the message types of the protocol; regenerate it rather than editing.
syntax = "proto3";

package place_order;

// What the client orders.
message CreateOrder {
  string item = 1;
  uint64 quantity = 2;
}

// The shop's price for the order, in cents.
message Quote {
  uint64 price = 1;
}

// The client pays; the card number is a secret, never logged.
message Pay {
  uint64 amount = 1;
  string card_secr = 2;
}

message Charge {
  uint64 amount = 1;
  string card_secr = 2;
}

message Receipt {
  uint64 amount = 1;
}

message Charged {
}

message Declined {
}

message PaymentDeclined {
}

message SoldOut {
}

// The event envelope: every event of the protocol, as it is carried and logged. Integers travel
// as decimal strings in JSON (the protobuf JSON mapping), so nanosecond timestamps stay exact.
message Event {
  uint64 ts = 1;        // nanoseconds; unique in its stream, and its identity
  string id = 2;        // the event's name, in a log that names events
  string session = 3;   // the session, or its key's value when the sessions block names one
  string from = 4;      // the sending actor: Role, or Role/instance
  string to = 5;        // the receiving actor
  string label = 6;
  string cause = 7;     // the id or ts of the event this one answers
  string prev = 8;      // the id or ts of the previous event the same actor sent
  uint64 sent_at = 9;   // the sender's clock when it sent, ns; compared only with that clock
  uint64 received_at = 10; // the receiver's clock when it received, ns
  string protocol = 11; // the protocol this event belongs to
  uint32 version = 12;  // its version; 0 when the protocol has none (fields 13 to 15 are kept for the header)
  oneof payload {
    CreateOrder create_order = 16;
    Quote quote = 17;
    Pay pay = 18;
    Charge charge = 19;
    Charged charged = 20;
    Receipt receipt = 21;
    Declined declined = 22;
    PaymentDeclined payment_declined = 23;
    SoldOut sold_out = 24;
  }
}
\end{lstlisting}


\section{The constructs}

The rest of the language adds repetition, concurrency and composition. Each construct
appears here with the smallest protocol that shows it; Chapter~\ref{ch:constructs}
gives each one's reduction and what it costs.

\paragraph{Loops.} \texttt{loop L \{ \ldots\ continue L; \}} repeats a body.
A loop that may end offers a branch without \texttt{continue}.

\begin{lstlisting}[language={},basicstyle=\ttfamily\footnotesize,breaklines=true]
protocol Poll (A, B) {
  rec L {
    choice {
      q: Ask from A to B;
      Answer from B to A cause q;
      continue L;
    } or {
      Bye from A to B;
    }
  }
}
\end{lstlisting}


\paragraph{Bounded loops.} \texttt{loop Try at most 3 \{ \ldots\ \} otherwise \{ \ldots\ \}}
repeats at most three times and then takes the \texttt{otherwise} block. LeanFM unrolls
it as soon as the text is read, so nothing later sees a counter; the counter becomes
three copies of the body. That makes bounded loops cheap to write and, as
Chapter~\ref{ch:constructs} shows, expensive to name in CTL.

\begin{lstlisting}[language={},basicstyle=\ttfamily\footnotesize,breaklines=true]
protocol Retry (A, B) {
  loop Try at most 3 {
    q: Ask from A to B;
    choice {
      Ok from B to A cause q;
    } or {
      Busy from B to A cause q;
      continue Try;
    }
  } otherwise {
    GiveUp from A to B;
  }
}
\end{lstlisting}


\paragraph{Parallel composition.} \texttt{parallel \{ \ldots\ \} and \{ \ldots\ \}} runs
branches whose messages interleave in any order. Each branch is a conversation of its
own; the protocol says nothing about how they interleave, so every interleaving is
allowed.

\begin{lstlisting}[language={},basicstyle=\ttfamily\footnotesize,breaklines=true]
protocol Both (A, B, C) {
  par {
    q: Ask from A to B;
    Answer from B to A cause q;
  } and {
    r: Ask2 from A to C;
    Answer2 from C to A cause r;
  }
}
\end{lstlisting}


\paragraph{Sessions.} \texttt{sessions \{ \ldots\ \}} marks a block that may run any
number of times beside the rest, each run a session of its own: a lobby where any
number of games are going on at once.

\section{Including one protocol in another}\label{sec:include}

\texttt{include} is how a protocol is assembled from smaller ones. It means exactly
\emph{paste the other protocol's body here}, after two substitutions: its roles are
renamed to roles of the includer, and the fields it was promised are renamed to fields
the includer already has. Here a site lets a customer log in, place an order, and log
out; placing an order is a protocol of its own. Two field types are new:
\texttt{Fresh} marks a value minted by that message (\texttt{LoggedIn} mints the session
id \texttt{sid}), and \texttt{Link} a value that names one minted earlier
(Chapter~\ref{ch:where} gives identity its own sections).

\begin{lstlisting}[language={},basicstyle=\ttfamily\footnotesize,breaklines=true]
protocol Order (Client, Shop) given token {
  c: Create from Client to Shop;
  Made from Shop to Client cause c;
}

protocol Site (Customer, Store) {
  l: Login from Customer to Store;
  LoggedIn from Store to Customer cause l;
  include Order Customer Store with token = sid;
  o: Logout from Customer to Store;
  Bye from Store to Customer cause o;
}
\end{lstlisting}


Read the two new pieces of syntax as a function definition and a call.
\begin{itemize}
\item \texttt{protocol Order (Client, Shop) given token} is the definition.
  \texttt{Client} and \texttt{Shop} are its parameters for roles. \texttt{given token}
  is a parameter for data: Order does not create the session id it runs under; whoever
  includes it must hand one in, and Order calls it \texttt{token}.
\item \texttt{include Order Customer Store with token = sid} is the call. The roles are
  passed by position, so Order's \texttt{Client} is played by \texttt{Customer} and its
  \texttt{Shop} by \texttt{Store}. \texttt{with token = sid} passes the data: the field
  Order calls \texttt{token} is the \texttt{sid} that \texttt{LoggedIn} minted. When
  the names already agree, \texttt{with sid} is short for \texttt{with sid = sid}.
\end{itemize}

So Site is the same protocol as this one, written out by hand, in which the order's
messages carry the field under the site's name \texttt{sid} instead of \texttt{token}
(the order's structures change accordingly):

\begin{lstlisting}[language={},basicstyle=\ttfamily\footnotesize,breaklines=true]
structure Login where
  user : String
structure LoggedIn where
  sid : Header (Fresh String)
structure Logout where
  sid : Header (Link String)
structure Create where
  sid : Header (Link String)
  item : String
structure Made where
  sid : Header (Link String)
  orderId : Fresh String

protocol Site (Customer, Store) {
  l: Login from Customer to Store;
  LoggedIn from Store to Customer cause l;
  c: Create from Customer to Store;
  Made from Store to Customer cause c;
  o: Logout from Customer to Store;
  Bye from Store to Customer cause o;
}
\end{lstlisting}


That is not an informal claim. LeanFM decides it by comparing the two (the faithfulness
check of Chapter~\ref{ch:faithful}):

\begin{lstlisting}[language={},basicstyle=\ttfamily\footnotesize,breaklines=true]
Site: the candidate (7 nodes) is faithful to the protocol (7 typestates), decided by walk
\end{lstlisting}


\begin{figure}[ht]
\centering
\AccessibleFigure{Interaction diagram of Site. Login from Customer to Store and LoggedIn back. Then a box labelled Order with token bound to sid containing Create from Customer to Store and Made back. Then Logout from Customer to Store and Bye back.}{\ThemedDiagram{width=.5\textwidth}{ctl-include}}
\caption{Site as LeanFM draws it: the included protocol is a box, labelled with how the id is passed in.}
\label{fig:include}
\end{figure}

Three things follow. Nothing in the reduction to CTL needs to know about
\texttt{include}: by the time a state graph is built, the included body is in place.
LeanFM checks the call like a type checker checks a function call: an included
protocol that is \texttt{given} a field must be passed it, the field passed must be in
scope and have the same type, and a field may be passed only once. And inside
\texttt{sessions \{ \ldots\ \}}, \texttt{include} is how one id becomes many sessions:
\texttt{sessions \{ include Order Customer Store with token = sid; \}} lets one login
place any number of orders, each its own session, each told apart by the
\texttt{orderId} its \texttt{Made} mints. GoServer's lobby includes its game this way,
and the game includes a seat (Chapter~\ref{ch:where}).

\section{The rest of the language}

\paragraph{Events.} \texttt{Finished at Server} is an event: something one role does
that no other role receives. Events order the conversation like messages do and
appear in CTL as \texttt{(on Finished)}.

\paragraph{What else a protocol may say.} A choice may name what drives it
(\texttt{choice at Server using declared("IsBusy")}) and give branch probabilities;
a message may say how long it takes (\texttt{after 50 ms}); a \texttt{require} block
may promise deadlines and fractions over observation periods; a \texttt{system} block
may run several sessions over shared, serial servers. These are LeanFM's extensions
for performance and deployment. They change the questions one can ask, not the
conversation's order, and this book, which is about order, sets them aside.

\section{A protocol is a set of runs}

Formally, a global type denotes the set of its \emph{runs}: the finite and infinite
sequences of messages and events it allows, each with the causes and field values
that make it a run of this protocol. The semantics LeanFM implements is operational.
A state is the part of the protocol still to be played, its \emph{residual}, together
with the ids already bound; a step plays one message the residual allows next; the
session may stop where the residual can be empty. Chapter~\ref{ch:checker} builds the
graph of these states and makes it the model CTL is evaluated over. Before that,
Chapter~\ref{ch:projection} asks whether the roles, each seeing only its own messages,
can play a protocol at all.
